\paragraph{Imported inputs.}
The proof uses the following results from Gubler--K\"unnemann
\cite{GK17}, with their original hypotheses.
\begin{enumerate}
\item Their \S9.7 constructs the global piecewise smooth part of the first
 Chern current and the current decomposition (9.7.1).
\item Their Proposition~9.8 constructs the residual generalized delta-form,
 proves that it has codimension one, and proves its uniqueness and
 functoriality.
\item Their Definition~9.9 gives the ambient-Hessian criterion for a
 delta-metric.  Their Remark~9.10 identifies its ambient superform with a
 generalized delta-preform of codimension zero.  Their Proposition~9.11 and
 \S9.13 identify the resulting sum with the unique first Chern delta-form.
\item Their Proposition~1.8 extends a piecewise smooth function from a
 tropical cycle to the ambient vector space.  Their Proposition~1.14 and
 Corollary~1.15 identify the corner locus after restriction with stable
 intersection of the ambient corner locus.  Their Corollary~3.19 is the
 tropical Poincar\'e--Lelong decomposition into facewise Hessian and corner
 locus.
\end{enumerate}
The bidegree convention is the one in \cite[Definition~2.9 and
Example~2.10]{GK17}: if \(D\) has ambient codimension \(l\), then
\(\alpha\wedge\delta_D\), with \(\alpha\) of bidegree \((p-l,q-l)\), has
ambient bidegree \((p,q)\).

\begin{proof}
\cref{itm:ch18-dm:frame}
The construction is \cite[\S9.7]{GK17}.  We spell out its invariance
because it prevents a misleading fixed-chart assertion.  If \(e'=ge\),
then
\[
  -\log\|e'\|=-\log\|e\|-\log|g|.
\]
After adjoining the invertible regular function \(g\) to a common very
affine refinement, \(-\log|g|\) is the restriction of an integral affine
function on the new ambient tropical vector space.  Its facewise Hessian
vanishes.  Hence the two local facewise Hessians agree on that common
refinement.  The same argument for two tropical charts, using their common
very-affine refinement, gives a global element of
\(\operatorname{PS}^{1,1}(W)\).  This construction takes place in the
piecewise smooth sheaf.  It supplies no ambient bidegree-\((1,1)\) lift, so
it does not by itself put the class in the codimension-zero part of the
generalized delta-form sheaf.

\cref{itm:ch18-dm:current}
Equation \eqref{eq:ch18-dm-current-decomposition} is
\cite[(9.7.1)]{GK17}.  Proposition~9.8 there constructs a generalized
delta-form of codimension one and proves exactly the displayed universal
pullback property.  It also proves uniqueness among generalized delta-forms
satisfying that property.  Taking \(f=\operatorname{id}_X\) gives the
residual summand in \eqref{eq:ch18-dm-current-decomposition}.  No injectivity
of the map from generalized delta-forms to associated currents on \(W\)
alone is used or asserted.

\cref{itm:ch18-dm:delta}
Condition \eqref{eq:ch18-dm-hessian-lift} is exactly
\cite[Definition~9.9(ii)]{GK17}; condition (i) of that definition is the
piecewise smooth local presentation already recorded in
\cref{itm:ch18-dm:frame}.  A superform with piecewise smooth coefficients on
the ambient vector space is a generalized delta-preform of codimension zero
by \cite[Remark~9.10 and Example~2.10]{GK17}.  That remark therefore proves
the stated equivalence.  Proposition~9.11 and \S9.13 of the same paper show
that the sum of this codimension-zero piece and the residual codimension-one
piece represents the full first Chern delta-current and is a delta-form.
Its uniqueness follows from \cite[Proposition~6.8 and
Definition~9.12]{GK17}.

\cref{itm:ch18-dm:degree}
By \cite[Proposition~1.8]{GK17}, after shrinking if necessary,
\(\phi_e\) has a piecewise smooth ambient extension
\(\widetilde\phi_e\).  The ambient corner locus
\(\widetilde\phi_e\cdot N_{U,\mathbf R}\) has codimension one.  Its stable
intersection with \(C\) is the corner locus of the restricted function:
\[
   (\widetilde\phi_e\cdot N_{U,\mathbf R})\cdot C
      =\phi_e\cdot C,
\]
by \cite[Proposition~1.14 and Corollary~1.15]{GK17}.  The facewise
restriction in \eqref{eq:ch18-dm-hessian-lift} gives the next identity on
all of \(C\), because \(C\subseteq\operatorname{trop}_U(V)\):
\[
   \gamma\wedge\delta_C
    =(\mathrm d'_{\mathrm P}\mathrm d''_{\mathrm P}\phi_e)
       \wedge\delta_C.
\]
Together with the tropical Poincar\'e--Lelong formula
\cite[Corollary~3.19]{GK17}, these two identities prove
\eqref{eq:ch18-dm-ambient-representative}.

It remains only to keep the two degree conventions separate.  The cycle
\(C\) has ambient codimension \(c\), so a bidegree-\((1,1)\) coefficient
form wedged with \(\delta_C\) has ambient bidegree
\((c+1,c+1)\).  The cycle
\((\widetilde\phi_e\cdot N_{U,\mathbf R})\cdot C\) has ambient codimension
\(c+1\), so its integration preform has the same ambient bidegree.  After
subtracting the ambient codimension \(c\) of \(C\), both have relative
bidegree \((1,1)\).  Their support codimensions are respectively \(c\) and
\(c+1\).  On the other hand \(\gamma\), before wedging with
\(\delta_C\), is by definition a codimension-zero generalized
delta-preform.  This proves every assertion in
\cref{itm:ch18-dm:degree}.
\end{proof}
